% Part: computability
% Chapter: computability-theory
% Section: computable-sets

\documentclass[../../../include/open-logic-section]{subfiles}

\begin{document}

\olfileid{cmp}{thy}{cps}
\olsection{مجموعه‌های محاسبه‌پذیر}

می‌توانیم مفهوم محاسبه‌پذیری را از توابع محاسبه‌پذیر به مجموعه‌های
محاسبه‌پذیر گسترش دهیم:

\begin{defn}
  بگذارید $S$ مجموعه‌ای از اعداد طبیعی باشد. آنگاه $S$
  \emph{محاسبه‌پذیر} است اگر و تنها اگر تابع مشخصهٔ آن~$\Char{S}$
  محاسبه‌پذیر باشد. به بیان دیگر، $S$ محاسبه‌پذیر است اگر و تنها اگر
  تابع
\[
\Char{S}(x) =
\begin{cases}
1 & \text{اگر $x \in S$} \\
0 & \text{در غیر این صورت}
\end{cases}
\]
محاسبه‌پذیر باشد. به همین ترتیب، رابطهٔ
$R(x_0,\dots,x_{k-1})$ محاسبه‌پذیر است اگر و تنها اگر تابع مشخصهٔ آن
محاسبه‌پذیر باشد.

مجموعه‌ها و روابط محاسبه‌پذیر را \emph{تصمیم‌پذیر} نیز می‌نامند.
\end{defn}

\begin{explain}
توجه کنید که اکنون چند مفهوم محاسبه‌پذیری داریم: برای توابع جزئی،
برای توابع و برای مجموعه‌ها. آن‌ها را با هم اشتباه نگیرید!{}
\iftag{TMs}{ماشین تورینگی که تابعی جزئی را محاسبه می‌کند، برای
  ورودی‌هایی که تابع در آن‌ها تعریف شده است خروجی تابع را بازمی‌گرداند؛
  ماشین تورینگی که مجموعه‌ای را محاسبه می‌کند، پس از تصمیم‌گیری دربارهٔ
  عضویت یا عدم عضویت مقدار ورودی در مجموعه، $1$ یا~$0$ را بازمی‌گرداند.}{}
\end{explain}

\end{document}
